\chapter{Integrations and Tool Contracts}

\section{Connector architecture}

Use three layers:

\begin{enumerate}
\item \textbf{Tool definition:} what the AI may request, with a stable name, description, input schema, output schema, and risk classification.
\item \textbf{Tool execution service:} validation, authorization, idempotency, audit, timeout, and result normalization.
\item \textbf{Connector adapter:} provider-specific authentication, endpoint calls, payload mapping, and error interpretation.
\end{enumerate}
Example: \code{AI engine \ensuremath{\rightarrow} orders.\allowbreak{}lookup tool \ensuremath{\rightarrow} tool execution service \ensuremath{\rightarrow} commerce connector \ensuremath{\rightarrow} commerce API}

The AI should never receive raw integration credentials or unrestricted HTTP access.

\section{Tool definition fields}

\begin{itemize}
\item stable tool name, e.g. \code{orders.\allowbreak{}lookup};
\item human-readable description;
\item version;
\item JSON input schema;
\item output schema;
\item allowed agent/tenant scope;
\item required permission;
\item risk level: read-only, reversible write, sensitive write;
\item confirmation/approval requirement;
\item timeout and retry policy;
\item idempotency behavior;
\item audit and redaction policy.
\end{itemize}
\section{Example: read-only order lookup}

\begin{lstlisting}
{
  "name": "orders.lookup",
  "version": 1,
  "risk": "read_only",
  "input_schema": {
    "type": "object",
    "properties": {
      "order_reference": {
        "type": "string",
        "minLength": 3,
        "maxLength": 80
      }
    },
    "required": ["order_reference"],
    "additionalProperties": false
  },
  "output_schema": {
    "type": "object",
    "properties": {
      "status": {"type": "string"},
      "estimated_delivery": {"type": ["string", "null"]}
    },
    "required": ["status", "estimated_delivery"],
    "additionalProperties": false
  }
}
\end{lstlisting}

This is a contract example, not a claim that any particular commerce provider uses this schema.

\section{Error mapping}

Normalize connector failures into categories:

\begin{itemize}
\item \code{AUTHENTICATION\_\allowbreak{}FAILED}
\item \code{AUTHORIZATION\_\allowbreak{}DENIED}
\item \code{NOT\_\allowbreak{}FOUND}
\item \code{VALIDATION\_\allowbreak{}FAILED}
\item \code{RATE\_\allowbreak{}LIMITED}
\item \code{TIMEOUT}
\item \code{PROVIDER\_\allowbreak{}UNAVAILABLE}
\item \code{CONFLICT}
\item \code{UNKNOWN\_\allowbreak{}OUTCOME}
\end{itemize}
Only retry failures that are safe to retry. For writes, an unknown outcome may require a status lookup or human review instead of blind retry.

\section{Webhook handling}

\begin{enumerate}
\item Verify provider signature and timestamp.
\item Resolve the configured channel/integration from trusted endpoint configuration.
\item Enforce body size and schema limits.
\item Persist event ID and payload/minimal normalized data.
\item Return success only after durable acceptance.
\item Process asynchronously.
\item Deduplicate provider retries.
\item Record processing state and failure reason without leaking secrets.
\end{enumerate}
\mmdfigL{epics/E06-safe-business-tools/02-three-layer-connector-architecture}{Three-layer connector architecture}{fig:siness-tools-02-three-layer-connector-architecture-166}


